\documentclass[11pt,a4paper]{article}
\usepackage[T1]{fontenc}
\usepackage[utf8]{inputenc}
\usepackage[main=italian,provide=*]{babel}
\usepackage{amsmath,amssymb}
\usepackage{geometry}
\geometry{margin=2.3cm,top=2.5cm,bottom=2.7cm}
\usepackage{booktabs}
\usepackage[table]{xcolor}
\usepackage{tcolorbox}
\tcbuselibrary{skins,breakable}
\usepackage{titlesec}
\usepackage{enumitem}
\usepackage{seqsplit}
\usepackage{needspace}
\usepackage{fancyhdr}
\usepackage{hyperref}
\hypersetup{colorlinks=true,linkcolor=Sepia,citecolor=Sepia,urlcolor=Sepia}

\definecolor{inkblue}{HTML}{1B3A5C}
\definecolor{lightblue}{HTML}{EAF1F8}
\definecolor{accent}{HTML}{B0562A}
\definecolor{softgray}{HTML}{6B6B6B}
\definecolor{okgreen}{HTML}{2E6B3E}
\definecolor{okgreenbg}{HTML}{EAF5EC}
\definecolor{warnamber}{HTML}{9C6B12}
\definecolor{warnbg}{HTML}{FBF1E0}
\definecolor{advisorpurple}{HTML}{5B3A7A}
\definecolor{advisorbg}{HTML}{F1ECF6}
\definecolor{notebar}{HTML}{9C6B12}

\titleformat{\section}
  {\normalfont\Large\bfseries\color{inkblue}}
  {\thesection}{0.6em}{}[{\color{inkblue!40}\titlerule}]
\titlespacing{\section}{0pt}{0.8em}{0.4em}
\titleformat{\subsection}
  {\normalfont\large\bfseries\color{accent}}
  {}{0em}{}
\titlespacing{\subsection}{0pt}{0.5em}{0.2em}

\pagestyle{fancy}
\fancyhf{}
\renewcommand{\headrulewidth}{0.4pt}
\fancyhead[R]{\small\color{softgray}\thepage}
\fancyfoot[C]{\small\color{softgray}Tesi triennale, Samuele Schiavi --- Relatore: Prof.\ Alessandro Chiesa}

\newtcolorbox{keyeq}[1][]{
  colback=lightblue, colframe=inkblue!70, boxrule=0.6pt,
  arc=2pt, left=8pt, right=8pt, top=4pt, bottom=4pt, #1
}
\newtcolorbox{panelbox}[1]{
  colback=white, colframe=softgray!50, boxrule=0.5pt,
  arc=2pt, left=7pt, right=7pt, top=3pt, bottom=3pt,
  fonttitle=\bfseries\color{inkblue}, title={#1}
}
\newtcolorbox{okbox}[1]{
  colback=okgreenbg, colframe=okgreen!60, boxrule=0.6pt,
  arc=2pt, left=7pt, right=7pt, top=3pt, bottom=3pt,
  fonttitle=\bfseries\color{okgreen}, title={#1}
}
\newtcolorbox{warnbox}[1]{
  colback=warnbg, colframe=warnamber!70, boxrule=0.6pt,
  arc=2pt, left=7pt, right=7pt, top=3pt, bottom=3pt,
  fonttitle=\bfseries\color{warnamber}, title={#1}
}
\newtcolorbox{advisorbox}[1]{
  colback=advisorbg, colframe=advisorpurple!60, boxrule=0.6pt,
  arc=2pt, left=7pt, right=7pt, top=4pt, bottom=4pt,
  fonttitle=\bfseries\color{advisorpurple}, title={#1}
}
\newtcolorbox{editnote}{
  colback=white, colframe=white, boxrule=0pt,
  borderline west={2.2pt}{0pt}{notebar},
  arc=0pt, left=10pt, right=4pt, top=2pt, bottom=2pt,
  fontupper=\itshape\small\color{softgray!90!black}
}
\fancyhead[L]{\small\color{softgray}Come si dividono i compiti: analisi\_rotazioni\_PMA vs confronto\_ansatz\_entangler}

\title{\vspace{-1.5em}\color{inkblue}Come si dividono i compiti\\[0.15em]
\Large\normalfont\color{softgray}\texttt{analisi\_rotazioni\_PMA.ipynb} e \texttt{confronto\_ansatz\_entangler.ipynb}}
\author{Note di lavoro --- Tesi triennale, Samuele Schiavi\\ Relatore: Prof.\ Alessandro Chiesa}
\date{}

\begin{document}
\sloppy
\maketitle
\thispagestyle{fancy}
\vspace{-2em}

\begin{center}
\renewcommand{\arraystretch}{1.5}
\small
\begin{tabular}{@{}p{2.0cm}p{5.7cm}p{5.7cm}@{}}
\toprule
\rowcolor{inkblue!10}
 & \textbf{\seqsplit{analisi\_rotazioni\_PMA.ipynb}} & \textbf{\seqsplit{confronto\_ansatz\_entangler.ipynb}}\\
\midrule
\textbf{Domanda} &
Dove mettere le rotazioni che rompono $M$: un qubit o due? Parametro
condiviso o indipendente? &
Come si comporta \emph{l'intera famiglia} di ansatz (PMA + 4 generici)
su tutto l'asse $b/J$ e al crescere di $D$?\\
\textbf{Metodo} &
4 esperimenti mirati (A--D) a un solo punto fisso, con derivazione
analitica del perché il parametro condiviso fallisce &
Sweep sistematico su griglia $(b,D)$, multistart, confronto con 11
ansatz insieme\\
\textbf{Esito} &
Stabilisce \emph{la regola}: serve un grado di libertà indipendente per
settore ($M=\pm1$); PMA-2q raggiunge $\mathcal F=1$ da 3 parametri, il
minimo teorico &
\emph{Applica} quella regola già decisa: usa
\texttt{PMA\_RECOMMENDED = "PMA-2q.3"} come riferimento canonico, e
mostra il quadro completo con anche gli ansatz generici A--D per
contrasto\\
\bottomrule
\end{tabular}
\end{center}

\begin{okbox}{In sintesi}
\texttt{analisi\_rotazioni\_PMA.ipynb} è il notebook di indagine mirata
che risolve una domanda di design puntuale (con derivazioni ed
esperimenti che \emph{falliscono} apposta, per isolare la causa).
\texttt{confronto\_ansatz\_entangler.ipynb} viene dopo, importa quella
conclusione già pronta (il codice stesso lo dimostra: le funzioni
\texttt{pma\_1q} e \texttt{pma\_2q} sono identiche, riga per riga, in
entrambi i notebook) e la mette a sistema nel confronto finale a 11
ansatz che stabilisce l'ansatz canonico per tutto il resto del
progetto.
\end{okbox}

\end{document}
